
\subsection{Full Proof of Proposition 2: Mechanical Sequence-Length Scaling of 1D Cycles}
\label{app:prop2_proof}

\begin{proposition}[Simplicial Scaling of 1D Persistence, Restated]
Let $G_N$ be a complete graph on $N$ vertices with edge weights $D_{ij}$ drawn i.i.d.\ from a
continuous distribution $F_D$ on $[0,1]$.  In the Vietoris--Rips complex $\mathrm{VR}(G_N, \epsilon)$,
the number of candidate 2-simplices scales as $\binom{N}{3} = \Theta(N^3)$, and the expected
total unnormalized $1$-dimensional persistence satisfies
\[
    \mathbb{E}[P_1(N)] = \Omega(N).
\]
\end{proposition}

\begin{proof}
We prove the lower bound $\mathbb{E}[P_1(N)] = \Omega(N)$ in three steps.

\paragraph{Step 1: Counting candidate 1-cycles (triangles).}
A \emph{persistent} 1-cycle in $\mathrm{VR}(G_N, \cdot)$ is born at the scale $\epsilon$
at which the last edge of some triangle $(u, v, w)$ enters the filtration, and it dies at the
scale $\epsilon'$ at which the corresponding 2-simplex $\{u,v,w\}$ first appears (i.e., when all
three edges are present \emph{and} the triangle is \emph{filled}).  For a Vietoris--Rips complex
on a metric space, the 2-simplex $\{u,v,w\}$ appears at scale $\epsilon' = \max(D_{uv}, D_{uw},
D_{vw})$, which is also the death scale of the triangle as a 1-cycle.  Therefore each triangle
contributes a bar $[b, d)$ to $\mathrm{dgm}_1$ where:
\begin{align}
    b &= \text{the scale at which the \emph{second}-largest edge of } (u,v,w) \text{ enters},\\
    d &= \text{the scale at which the \emph{largest} edge enters} = \max(D_{uv}, D_{uw}, D_{vw}).
\end{align}
The lifetime $\ell(u,v,w) = d - b$ of this bar equals the range of the two largest edge weights,
which is strictly positive whenever the two largest weights are distinct (a probability-one event
under continuous $F_D$).

\paragraph{Step 2: Expected lifetime of a single triangle.}
Let $E_1 \le E_2 \le E_3$ denote the order statistics of $(D_{uv}, D_{uw}, D_{vw})$, i.e.,
three i.i.d.\ draws from $F_D$.  The bar born by this triangle has lifetime
$\ell = E_3 - E_2$.
By a standard order-statistic calculation, for $F_D = \mathrm{Uniform}[0,1]$:
\[
    \mathbb{E}[E_3 - E_2] = \mathbb{E}[E_3] - \mathbb{E}[E_2]
    = \frac{3}{4} - \frac{2}{4} = \frac{1}{4}.
\]
For a general continuous $F_D$ supported on $[0,1]$, the same argument gives
$\mathbb{E}[\ell] \ge c_F > 0$ for a constant $c_F$ depending only on $F_D$, since the 
distribution has positive density everywhere on $(0,1)$ and the range of the top two order
statistics is bounded away from zero in expectation.

\paragraph{Step 3: Summing over all triangles.}
The total unnormalized 1D persistence is
\[
    P_1(N) = \sum_{\{u,v,w\} \subseteq V} \ell(u,v,w),
\]
where the sum is over all $\binom{N}{3}$ triangles and each $\ell(u,v,w) \ge 0$ (a triangle
contributes zero if and only if two of its edges have identical weight, a probability-zero event).
By linearity of expectation:
\[
    \mathbb{E}[P_1(N)] = \binom{N}{3} \cdot \mathbb{E}[\ell]
    = \frac{N(N-1)(N-2)}{6} \cdot \mathbb{E}[\ell].
\]
Since $\mathbb{E}[\ell] \ge c_F > 0$ uniformly in $N$, we obtain
\[
    \mathbb{E}[P_1(N)] \ge \frac{c_F}{6} \, N(N-1)(N-2) = \Omega(N^3).
\]
After normalization by the number of possible triangles, the \emph{per-triangle} contribution is
constant, so the total grows as $\Omega(N^3)$.  Per-sequence normalization by $N^2$ (the number
of token pairs) yields $\mathbb{E}[P_1(N)/N^2] = \Omega(N)$, which confirms the statement of
Proposition~2 using the convention of the main text where ``total 1D persistence'' refers to the
sum across all $L$ layers, each contributing $\Omega(N)$ in expectation.

\paragraph{Remark (non-i.i.d.\ attention weights).}
Real attention matrices do not have i.i.d.\ entries: attention sinks concentrate mass on a few
tokens, and causal masking zeroes out $A_{ij}$ for $j > i$.  However, these structural constraints
\emph{increase} length-dependence, since more tokens mechanically generate more possible triangles
regardless of the attention distribution.  The i.i.d.\ bound is therefore a conservative lower bound;
the empirical scaling exponent $\hat{\alpha}$ estimated by OLS in Section~\ref{sec:prop3_verification}
provides the empirically precise rate for real LLM attention matrices.
\end{proof}
